\begin{definition}[Injective inverse map for a simple MPO tensor]
    \label{def:mpdo_injective_inverse_layer}
    \lean{MPOTensor.IsInjective}
    \lean{MPOTensor.inverseTensor}
    \lean{MPOTensor.physRealize}
    \lean{MPOTensor.physRealizeLeft}
    \leanok
    \uses{def:mpo_tensor, def:mpo_to_mps, def:injective}
    A simple MPO tensor $K$ is called injective if its doubled-index MPS tensor
    is injective. For such a tensor, the chosen decomposition map of the
    doubled-index MPS furnishes a concrete inverse tensor $K^{-1}$ together
    with right- and left-physical realization maps that turn virtual bond
    insertions back into local physical operators.
\end{definition}

\begin{theorem}[Inverse-map identities for an injective simple MPO tensor]
    \label{thm:mpdo_injective_inverse_layer_spec}
    \lean{MPOTensor.inverseTensor_spec}
    \lean{MPOTensor.physRealize_spec}
    \lean{MPOTensor.physRealize_mul}
    \lean{MPOTensor.physRealizeLeft_spec}
    \leanok
    \uses{def:mpdo_injective_inverse_layer}
    For an injective simple MPO tensor, contracting $K^{-1}$ with $K$ recovers
    the matrix units on the virtual bond space. Likewise, every right virtual
    insertion admits a physical realization, this realization is multiplicative,
    and the analogous left-insertion statement also holds.
\end{theorem}

For an injective simple tensor, contracting $K^{-1}$ with $K$ through the
physical index identifies the corresponding virtual legs. In components,
this is the matrix-unit identity
\begin{align}
    \sum_p(K^{-1})^{\alpha,\beta}_p(K^p)_{\alpha',\beta'}
    &=\delta_{\alpha,\alpha'}\delta_{\beta,\beta'}.
    \label{eq:rfp_inverse_matrix_units}
\end{align}
Its two Kronecker factors are the two bare cups:
\begin{tenkzequation}
    % The upper inverse tensor and lower tensor contract through the doubled
    % physical index p.  Their four virtual legs remain open and are
    % identified pairwise by the tensor product of the two cups on the right.
    \tenkzkernel
    \begin{tenkz}[rows={ket,bra}, cols=1]
        \tn[skin=box, ports={180:virtual, 0:virtual}]{K^{-1}} \\
        \tn[skin=box, ports={180:virtual, 0:virtual}]{K}
    \end{tenkz}
    $ = $
    \begin{tenkz}[rows={wire,wire}, cols=1, west=open, east=cup, bonds=none]
        \tn[skin=none]{} \\ \tn[skin=none]{}
    \end{tenkz}
    $ \otimes $
    \begin{tenkz}[rows={wire,wire}, cols=1, west=cup, east=open, bonds=none]
        \tn[skin=none]{} \\ \tn[skin=none]{}
    \end{tenkz}
\end{tenkzequation}
This is the identity in
\cite[Appendix~C.2, lines~1376--1380]{Cirac2016MPDO_arXiv}.

Applying this inverse map to the doubled physical index and closing its
virtual legs against a further copy of $K$ recovers $K$:
\begin{tenkzequation}
    % K K^{-1} K = K (arXiv:1606.00608, eq:Kidentity): the top
    % tensor closes both virtual legs against K^{-1} and keeps its physical
    % leg p open; K^{-1} contracts through the doubled physical index only
    % with the bottom K, whose virtual legs remain open.
    % Both sides have boundary signature (virtual-west=1 | virtual-east=1 |
    % physical-up=1 | physical-down=0).
    \tenkzkernel
    \begin{tenkz}[rows={ket,ket,bra}, cols=1, west=cup, east=cup, bonds=none]
        \tn[skin=box, ports={90:physical:$p$}]{K} \\
        \tn[skin=box, name=Kinv, ports={270:physical}]{K^{-1}} \\
        \tn[skin=box, name=Kbot,
            ports={90:physical, 180:virtual, 0:virtual}]{K}
        \tnwire{Kinv.270}{Kbot.90}
    \end{tenkz}
    $ = $
    \begin{tenkz}[cols=1, west=open, east=open]
        \tn[skin=box, ports={90:physical:$p$}]{K}
    \end{tenkz}
\end{tenkzequation}
This is the single-block identity in
\cite[Appendix~C.2, lines~1671--1676]{Cirac2016MPDO_arXiv}.

\begin{definition}[Normalized fourth-site virtual tail]
    \label{def:mpdo_normalized_four_site_tail}
    \lean{MPOTensor.normalizedFourSiteTail}
    \leanok
    \uses{def:mpo_family, def:mpo_physical_trace_transfer}
    Let $\mathcal K$ be an MPO tensor and set
    $Z_4=\tr[\rho^{(4)}(\mathcal K)]$. Its normalized fourth-site virtual tail is
    \begin{align}
        R_4
        &:=Z_4^{-1}\mathcal T_{\mathcal K}
          =\tr[\rho^{(4)}(\mathcal K)]^{-1}
          \sum_{i_4=0}^{d-1}\mathcal K^{i_4i_4}.
        \label{eq:rfp_normalized_tail}
    \end{align}
    Thus the contraction of the fourth ket and bra indices is incorporated into
    the virtual closing matrix. This is the $N=4$ specialization of the
    normalization convention and three-site marginal in
    \cite[lines~792--793 and 1343--1348]{Cirac2016MPDO_arXiv}.
\end{definition}

\begin{theorem}[Three-site marginal closed by the normalized fourth-site tail]
    \label{thm:mpdo_reduced_four_three_tail}
    \lean{MPOTensor.reducedBlockState_four_three_apply}
    \leanok
    \uses{def:mpdo_normalized_four_site_tail, def:mpo_eval_word}
    Let $u=(i_1,i_2,i_3)$ and $v=(j_1,j_2,j_3)$ be physical words of length
    three. The corresponding entry of the first-three-site marginal of the
    normalized four-site periodic operator is
    \begin{align}
        \sigma_3^{(4)}(\mathcal K)_{u,v}
        &=\tr\!\left(\mathcal K^{u,v}R_4\right)
          =\tr\!\left(\mathcal K^{i_1j_1}\mathcal K^{i_2j_2}
          \mathcal K^{i_3j_3}R_4\right).
        \label{eq:rfp_three_site_marginal}
    \end{align}
\end{theorem}

\begin{figure}[ht]
    \centering
    % Tracing the fourth physical index folds that site into R_4 and leaves
    % the first three ket and bra legs open.  The periodic wire closes the
    % virtual ring through the virtual-only R_4 cell.  The boundary signature
    % is (virtual-west=0 | virtual-east=0 |
    % physical-up=3 | physical-down=3).
    \tenkzkernel
    \begin{tenkz}[cols=4, boundary=periodic]
        \tn[label pos=w, ports={90:physical:$i_1$, 270:physical:$j_1$}]{\mathcal K} &
        \tn[label pos=w, ports={90:physical:$i_2$, 270:physical:$j_2$}]{\mathcal K} &
        \tn[label pos=w, ports={90:physical:$i_3$, 270:physical:$j_3$}]{\mathcal K} &
        \tn[skin=ring]{R_4}
    \end{tenkz}
    \caption{Tracing the fourth site of the normalized four-site periodic MPO
    leaves the physical legs of sites $1,2,3$ open and replaces the fourth site
    by the virtual closing operator
    $R_4=\tr[\rho^{(4)}(\mathcal K)]^{-1}
    \sum_{i_4}\mathcal K^{i_4 i_4}$; compare the normalization convention and
    three-site marginal in
    \cite[lines~792--793 and 1343--1348]{Cirac2016MPDO_arXiv}.}
    \label{fig:mpdo_normalized_four_site_tail}
\end{figure}

\begin{proof}
    \leanok
    \uses{def:mpdo_normalized_four_site_tail,
        lem:mpo_eval_word_append}
    Expanding the partial trace over the fourth site and factoring each
    length-four word at its last letter gives
    \begin{align}
        \sigma_3^{(4)}(\mathcal K)_{u,v}
        &=Z_4^{-1}\sum_{i_4=0}^{d-1}
          \tr\!\left(\mathcal K^{u,v}\mathcal K^{i_4i_4}\right).
        \notag
    \end{align}
    Linearity of the trace and (\ref{eq:rfp_normalized_tail}) now yield
    \begin{align}
        Z_4^{-1}\sum_{i_4=0}^{d-1}
          \tr\!\left(\mathcal K^{u,v}\mathcal K^{i_4i_4}\right)
        &=\tr\!\left(\mathcal K^{u,v}
          \left(Z_4^{-1}\sum_{i_4=0}^{d-1}\mathcal K^{i_4i_4}\right)\right)
          =\tr\!\left(\mathcal K^{u,v}R_4\right).
        \notag
    \end{align}
\end{proof}

\begin{theorem}[Non-vanishing of the normalized fourth-site tail]
    \label{thm:mpdo_normalized_four_site_tail_ne_zero}
    \lean{MPOTensor.normalizedFourSiteTail_ne_zero}
    \leanok
    \uses{def:mpdo_normalized_four_site_tail}
    If $Z_4=\tr[\rho^{(4)}(\mathcal K)]\ne0$, then $R_4\ne0$.
\end{theorem}

\begin{proof}
    \leanok
    \uses{thm:mpdo_reduced_four_three_tail}
    If $R_4=0$, then the tail formula makes every entry of
    $\sigma_3^{(4)}(\mathcal K)$ vanish, so
    $\tr[\sigma_3^{(4)}(\mathcal K)]=0$. On the other hand, $Z_4\ne0$ implies
    $\tr[\sigma_3^{(4)}(\mathcal K)]=1$, a contradiction.
\end{proof}

\begin{theorem}[A nonzero entry of the normalized fourth-site tail]
    \label{thm:mpdo_normalized_four_site_tail_entry_ne_zero}
    \lean{MPOTensor.exists_normalizedFourSiteTail_entry_ne_zero}
    \leanok
    \uses{def:mpdo_normalized_four_site_tail}
    If $Z_4\ne0$, there exist virtual indices $\beta$ and $\alpha$ such that
    $(R_4)_{\beta,\alpha}\ne0$. These are the indices selected before the
    sector factorization in
    \cite[Appendix~C.2, lines~1431--1437]{Cirac2016MPDO_arXiv}.
\end{theorem}

\begin{proof}
    \leanok
    \uses{thm:mpdo_normalized_four_site_tail_ne_zero}
    If every matrix entry of $R_4$ vanished, then $R_4$ would be the zero
    matrix, contrary to the preceding theorem.
\end{proof}

\begin{definition}[Three-site inverse-map contraction]
    \label{def:mpdo_inverse_map_three_site_contraction}
    \lean{MPOTensor.inverseMapThreeSiteContraction}
    \leanok
    \uses{def:mpdo_injective_inverse_layer}
    Let $\mathcal K$ be an injective simple MPO tensor, and let $R\in\MN{D}$ be
    the virtual contraction of the sites following a three-site word. For
    virtual indices
    $\alpha_1,\beta_1,\alpha_3,\beta_3$, define
    \begin{align}
        C_{\alpha_1,\beta_1,\alpha_3,\beta_3}(p_2;R)
        &:=
        \sum_{p_1,p_3}
        (\mathcal K^{-1})^{\alpha_1,\beta_1}_{p_1}
        (\mathcal K^{-1})^{\alpha_3,\beta_3}_{p_3}
        \tr\!\left(\mathcal K^{p_1}\mathcal K^{p_2}\mathcal K^{p_3}R\right).
        \label{eq:rfp_inverse_contraction}
    \end{align}
    This is the contraction used in the proof of
    \cite[Appendix~C.2, lines~1415--1438]{Cirac2016MPDO_arXiv}.
\end{definition}

\begin{theorem}[Evaluation of the three-site inverse-map contraction]
    \label{thm:mpdo_inverse_map_three_site_contraction}
    \lean{MPOTensor.inverseMapThreeSiteContraction_eq}
    \leanok
    \uses{def:mpdo_inverse_map_three_site_contraction}
    % The source display following Qketc repeats alpha_3 in the tail scalar.
    % Direct matrix-unit contraction gives R_{beta_3,alpha_1}; see the paper-gap note.
    For every middle physical index $p_2$,
    \begin{align}
        C_{\alpha_1,\beta_1,\alpha_3,\beta_3}(p_2;R)
        &=\mathcal K^{p_2}_{\beta_1,\alpha_3}
          R_{\beta_3,\alpha_1}.
        \label{eq:rfp_inverse_contraction_eval}
    \end{align}
    % Local fix (tail index): the source display repeats the third-site left
    % index in the tail entry; direct contraction gives R_{beta_3,alpha_1} as
    % above. Recorded in
    % docs/paper-gaps/cpgsv17_mpdo_sal_zcl_eta_local_structure.tex.
\end{theorem}

\begin{figure}[ht]
    \centering
        % (\mathcal K^{-1} (x) Id (x) \mathcal K^{-1}) \mathcal K_3
        %   = \mathcal K^{p_2}_{beta_1,alpha_3} R_{beta_3,alpha_1}
        % (arXiv:1606.00608, Appendix C.2, lines 1415--1438; paper
        % figures MPDO_Xab.png, MPDO_kappaab.png, and MPDO_mba.png).
        %
        % LHS: the two \mathcal K^{-1} boxes expose virtual pairs
        % (alpha_1,beta_1) and (alpha_3,beta_3).  Their physical indices
        % p_1,p_3 match the first and third legs of the three-site chain
        % closed against R; p_2 remains open.  RHS: \mathcal K carries beta_1 west,
        % p_2 up, and alpha_3 east, while R carries beta_3 west and
        % alpha_1 east.  There is no surviving sum: the identity is
        % pointwise in all four virtual indices.  Both sides have aggregate
        % boundary signature (virtual-west=2 | physical-up=1 |
        % virtual-east=2).
        \begin{tenkzequation}
            \tenkzkernel
            $\Bigl($
            \begin{tenkz}[cols=1, size=s]
                \tn[label pos=n, ports={180:virtual:$\alpha_1$,
                    0:virtual:$\beta_1$, 270:physical:$p_1$}]{\mathcal K^{-1}}
            \end{tenkz}%
            $\;\otimes\;\Id\;\otimes\;$%
            \begin{tenkz}[cols=1, size=s]
                \tn[label pos=n, ports={180:virtual:$\alpha_3$,
                    0:virtual:$\beta_3$, 270:physical:$p_3$}]{\mathcal K^{-1}}
            \end{tenkz}
            $\Bigr)$
            \begin{tenkz}[cols=4, boundary=periodic, size=s]
                \tn[label pos=s, ports={90:physical:$p_1$}]{\mathcal K} &
                \tn[label pos=s, ports={90:physical:$p_2$}]{\mathcal K} &
                \tn[label pos=s, ports={90:physical:$p_3$}]{\mathcal K} &
                \tn[skin=ring]{R}
            \end{tenkz}
            $ = $
            \begin{tenkz}[cols=1, size=s]
                \tn[label pos=s, ports={180:virtual:$\beta_1$,
                    0:virtual:$\alpha_3$, 90:physical:$p_2$}]{\mathcal K}
            \end{tenkz}
            $ \otimes $
            \begin{tenkz}[cols=1, size=s]
                \tn[skin=ring, ports={180:virtual:$\beta_3$,
                    0:virtual:$\alpha_1$}]{R}
            \end{tenkz}
        \end{tenkzequation}
    \caption{Applying $\mathcal K^{-1}$ at the first and third sites collapses
    the three-site contraction to
    $\mathcal K^{p_2}_{\beta_1,\alpha_3}R_{\beta_3,\alpha_1}$, where
    $p_2=(i_2,j_2)$ is the doubled physical index; compare
    \cite[Appendix~C.2, lines~1415--1438]{Cirac2016MPDO_arXiv}.}
    \label{fig:mpdo_inverse_map_three_site_contraction}
\end{figure}

\begin{proof}\leanok
    \uses{thm:mpdo_injective_inverse_layer_spec}
    Move the two finite sums through the trace. By the inverse-map identity
    \begin{align}
        \sum_p(\mathcal K^{-1})^{\alpha,\beta}_p\,\mathcal K^p
        &=E_{\alpha,\beta},
        \notag
    \end{align}
    which is the matrix form of (\ref{eq:rfp_inverse_matrix_units}), the double
    sum collapses to
    \begin{align}
        \tr\!\left(E_{\alpha_1,\beta_1}\mathcal K^{p_2}
          E_{\alpha_3,\beta_3}R\right)
        &=\mathcal K^{p_2}_{\beta_1,\alpha_3}
          R_{\beta_3,\alpha_1}.
        \notag
    \end{align}
\end{proof}

\begin{definition}[Three-site closure against a virtual tail]
    \label{def:mpdo_three_site_closure}
    \lean{MPOTensor.IsThreeSiteClosure}
    \leanok
    \uses{def:mpo_tensor}
    Let $\mathcal K$ be an MPO tensor and let $R\in\MN{D}$. A tripartite
    operator $\rho$ is the three-site closure of $\mathcal K$ against $R$ if
    \begin{align}
        \rho_{i_1i_2i_3,j_1j_2j_3}
        &=\tr\!\left(\mathcal K^{i_1j_1}\mathcal K^{i_2j_2}
          \mathcal K^{i_3j_3}R\right).
        \label{eq:rfp_three_site_closure}
    \end{align}
    This is the algebraic form of the three-site reduced operator in
    \cite[Appendix~C.2, lines~1422--1433]{Cirac2016MPDO_arXiv}.
\end{definition}

\begin{definition}[Physical slice and Hayashi inverse factors]
    \label{def:mpdo_hayashi_inverse_factors}
    \lean{MPOTensor.physicalSlice}
    \lean{MPOTensor.hayashiInverseLeft}
    \lean{MPOTensor.hayashiInverseRight}
    \leanok
    \uses{def:mpdo_injective_inverse_layer, def:mpdo_eta_structure}
    Fix virtual indices $\beta,\alpha$ and set
    $\kappa_{\beta,\alpha}(i,j)=\mathcal K^{ij}_{\beta,\alpha}$.
    For a Hayashi sector $k$, define the two inverse factors
    \begin{align}
        A^{(k)}_{\alpha_1,\beta_1}(l,l')
        &:=
        \sum_{i_1,j_1}
        (\mathcal K^{-1})^{\alpha_1,\beta_1}_{(i_1,j_1)}
        (\rho_{A b_1}^{(k)})_{(i_1,l),(j_1,l')},
        \notag\\
        B^{(k)}_{\alpha_3,\beta_3}(r,r')
        &:=
        \sum_{i_3,j_3}
        (\mathcal K^{-1})^{\alpha_3,\beta_3}_{(i_3,j_3)}
        (\rho_{b_2 C}^{(k)})_{(r,i_3),(r',j_3)}.
        \notag
    \end{align}
    These are the two factors in the sectorwise inverse-map calculation of
    \cite[Appendix~C.2, lines~1407--1428]{Cirac2016MPDO_arXiv}.
\end{definition}

\begin{theorem}[Outer inverse-map contraction of a closure]
    \label{thm:mpdo_closure_outer_contraction}
    \lean{MPOTensor.inverseMap_threeSite_closure_collapse}
    \lean{MPOTensor.inverseMap_conj_physicalSlice_expansion}
    \leanok
    \uses{def:mpdo_three_site_closure, def:mpdo_hayashi_inverse_factors}
    Let $\mathcal K$ be injective and let $\rho$ be its three-site closure
    against $R$. Applying the inverse tensor at the two outer sites collapses
    the closure to one entry of the middle physical slice and one entry of the
    tail:
    \begin{align}
        &\sum_{i_1,j_1,i_3,j_3}
        (\mathcal K^{-1})^{\alpha_1,\beta_1}_{(i_1,j_1)}
        (\mathcal K^{-1})^{\alpha_3,\beta_3}_{(i_3,j_3)}
        \rho_{i_1i_2i_3,j_1j_2j_3}
        \notag\\
        &\qquad
        =\mathcal K^{i_2j_2}_{\beta_1,\alpha_3}
        R_{\beta_3,\alpha_1}.
        \label{eq:rfp_outer_collapse}
    \end{align}
    Consequently, for any matrix $U$ on the middle site,
    \begin{align}
        &R_{\beta_3,\alpha_1}
        (U\kappa_{\beta_1,\alpha_3}U^\dagger)_{b,b'}
        \notag\\
        &\qquad
        =\sum_{i_1,j_1,i_3,j_3}
        (\mathcal K^{-1})^{\alpha_1,\beta_1}_{(i_1,j_1)}
        (\mathcal K^{-1})^{\alpha_3,\beta_3}_{(i_3,j_3)}
        \sum_{i_2,j_2}
        U_{b,i_2} \rho_{i_1i_2i_3,j_1j_2j_3}
        \overline{U_{b',j_2}}.
        \label{eq:rfp_conjugated_slice}
    \end{align}
\end{theorem}

\begin{proof}
    \leanok
    \uses{thm:mpdo_inverse_map_three_site_contraction}
    Substituting the three-site closure relation~
    (\ref{eq:rfp_three_site_closure}) converts the left-hand side
    of (\ref{eq:rfp_outer_collapse}) to
    $C_{\alpha_1,\beta_1,\alpha_3,\beta_3}((i_2,j_2);R)$.
    Theorem~\ref{thm:mpdo_inverse_map_three_site_contraction} then gives
    \begin{align}
        C_{\alpha_1,\beta_1,\alpha_3,\beta_3}((i_2,j_2);R)
        &=\mathcal K^{i_2j_2}_{\beta_1,\alpha_3}
          R_{\beta_3,\alpha_1},
        \notag
    \end{align}
    which is (\ref{eq:rfp_outer_collapse}). The second identity follows by
    expanding the matrix product entrywise,
    \begin{align}
        (U\kappa_{\beta_1,\alpha_3}U^\dagger)_{b,b'}
        &=\sum_{i_2,j_2}
        U_{b,i_2} \kappa_{\beta_1,\alpha_3}(i_2,j_2)
        \overline{U_{b',j_2}},
        \notag
    \end{align}
    multiplying by $R_{\beta_3,\alpha_1}$, substituting
    \begin{align}
        R_{\beta_3,\alpha_1}
        \kappa_{\beta_1,\alpha_3}(i_2,j_2)
        &=\sum_{i_1,j_1,i_3,j_3}
        (\mathcal K^{-1})^{\alpha_1,\beta_1}_{(i_1,j_1)}
        (\mathcal K^{-1})^{\alpha_3,\beta_3}_{(i_3,j_3)}
        \rho_{i_1i_2i_3,j_1j_2j_3},
        \notag
    \end{align}
    and reordering the finite sums.
\end{proof}

\begin{theorem}[Sectorwise inverse-map comparison]
    \label{thm:mpdo_inverse_map_hayashi_sector_comparison}
    \lean{MPOTensor.inverseMap_hayashi_sector_comparison}
    \leanok
    \uses{def:mpdo_three_site_closure,
        def:mpdo_hayashi_inverse_factors}
    Suppose that $\mathcal K$ is injective, $\rho$ is its three-site closure
    against $R$, and a Hayashi decomposition of $\rho$ has been chosen. Write
    \begin{align}
        \widetilde\kappa^{(k)}_{\beta_1,\alpha_3}
          ((l,r),(l',r'))
        &:=
        \left[U_B\kappa_{\beta_1,\alpha_3}U_B^\dagger\right]_{(k,l,r),(k,l',r')}.
        \notag
    \end{align}
    Then every diagonal sector satisfies
    \begin{align}
        R_{\beta_3,\alpha_1}
        \widetilde\kappa^{(k)}_{\beta_1,\alpha_3}
          ((l,r),(l',r'))
        &=p_k A^{(k)}_{\alpha_1,\beta_1}(l,l')
          B^{(k)}_{\alpha_3,\beta_3}(r,r').
        \label{eq:rfp_sector_comparison}
    \end{align}
    The probability $p_k$ remains explicit because both sector density
    operators are normalized. No zero-correlation-length or non-vanishing
    hypothesis is used in this comparison.
\end{theorem}

For fixed $k,\alpha_1,\beta_3$, the identity
(\ref{eq:rfp_sector_comparison}) reads coefficientwise as follows.
% Formula:
% R_{\beta_3,\alpha_1}
% \widetilde\kappa^{(k)}_{\beta_1,\alpha_3}((l,r),(l',r'))
% = p_k A^{(k)}_{\alpha_1,\beta_1}(l,l')
% B^{(k)}_{\alpha_3,\beta_3}(r,r').
% The displayed tensor has free indices beta_1 west, alpha_3 east, (l,r)
% north, and (l',r') south.  None is contracted, so its boundary signature is
% (1,1,1,1).  Its square is the coefficient tensor
% \widetilde\kappa^{(k)}; the four incident strokes carry precisely those four
% free indices.  The fixed k, alpha_1, and beta_3 select the sector and scalar
% coefficients.  R, A, and B remain scalar factors under ordinary
% multiplication; they are not tensor nodes.
% Source comparison: Papers/1606.00608/MPDO_kappaab.png, from lines 1415--1438
% of MPDO-22-12-17-2.tex, fixes the one-tensor four-leg topology.  The
% normalized formula is Theorem thm:mpdo_inverse_map_hayashi_sector_comparison.
\begin{tenkzequation}
\tenkzkernel
$R_{\beta_3,\alpha_1}\,$
\begin{tenkz}[cols=1]
    \tn[skin=mpo, ports={90:physical:$(l{,}r)$, 270:physical:$(l'{,}r')$,
        180:virtual:$\beta_1$, 0:virtual:$\alpha_3$}]{\widetilde\kappa^{(k)}}
\end{tenkz}
$ = $
$p_k\,A^{(k)}_{\alpha_1,\beta_1}(l,l')
B^{(k)}_{\alpha_3,\beta_3}(r,r')$.
\end{tenkzequation}

\begin{proof}
    \leanok
    \uses{thm:mpdo_closure_outer_contraction}
    At fixed indices $(k,l,r)$ and $(k,l',r')$, the Hayashi block equality is
    \begin{align}
        &\sum_{i_2,j_2}
        [U_B]_{(k,l,r),i_2}
        \rho_{i_1i_2i_3,j_1j_2j_3}
        \overline{[U_B]_{(k,l',r'),j_2}}
        \notag\\
        &\qquad
        =p_k(\rho_{A b_1}^{(k)})_{(i_1,l),(j_1,l')}
          (\rho_{b_2 C}^{(k)})_{(r,i_3),(r',j_3)}.
        \notag
    \end{align}
    For each $i_2,j_2$, the preceding theorem gives
    \begin{align}
        &\sum_{i_1,j_1,i_3,j_3}
        (\mathcal K^{-1})^{\alpha_1,\beta_1}_{(i_1,j_1)}
        (\mathcal K^{-1})^{\alpha_3,\beta_3}_{(i_3,j_3)}
        \rho_{i_1i_2i_3,j_1j_2j_3}
        \notag\\
        &\qquad
        =\mathcal K^{i_2j_2}_{\beta_1,\alpha_3}
          R_{\beta_3,\alpha_1}.
        \notag
    \end{align}
    Apply
    \begin{align}
        \sum_{i_1,j_1,i_3,j_3}
        &(\mathcal K^{-1})^{\alpha_1,\beta_1}_{(i_1,j_1)}
        (\mathcal K^{-1})^{\alpha_3,\beta_3}_{(i_3,j_3)}\,(-)
        \notag
    \end{align}
    to the Hayashi block equality. Substituting the preceding display and
    interchanging the finite sums gives
    \begin{align}
        &\sum_{i_1,j_1,i_3,j_3}
          (\mathcal K^{-1})^{\alpha_1,\beta_1}_{(i_1,j_1)}
          (\mathcal K^{-1})^{\alpha_3,\beta_3}_{(i_3,j_3)}
        \notag\\
        &\qquad{}\times\left(\sum_{i_2,j_2}
          [U_B]_{(k,l,r),i_2}
          \rho_{i_1i_2i_3,j_1j_2j_3}
          \overline{[U_B]_{(k,l',r'),j_2}}\right)
        \notag\\
        &\qquad
        =\sum_{i_2,j_2}
          [U_B]_{(k,l,r),i_2}
          \mathcal K^{i_2j_2}_{\beta_1,\alpha_3}
          R_{\beta_3,\alpha_1}
          \overline{[U_B]_{(k,l',r'),j_2}}.
        \notag
    \end{align}
    The right-hand side is
    \begin{align}
        &R_{\beta_3,\alpha_1}
        \sum_{i_2,j_2}
        [U_B]_{(k,l,r),i_2}
        \kappa_{\beta_1,\alpha_3}(i_2,j_2)
        \overline{[U_B]_{(k,l',r'),j_2}}
        \notag\\
        &\qquad
        =R_{\beta_3,\alpha_1}
        \widetilde\kappa^{(k)}_{\beta_1,\alpha_3}
          ((l,r),(l',r')).
        \notag
    \end{align}
    The contraction of the right-hand side of the Hayashi block equality is
    $p_k A^{(k)}_{\alpha_1,\beta_1}(l,l')
    B^{(k)}_{\alpha_3,\beta_3}(r,r')$.
    Since the Hayashi block equality holds, equating the two contractions
    gives (\ref{eq:rfp_sector_comparison}).
\end{proof}

\begin{theorem}[The four-site marginal is a three-site closure]
    \label{thm:mpdo_three_site_closure_reduced}
    \lean{MPOTensor.isThreeSiteClosure_reducedBlockState}
    \leanok
    \uses{def:mpdo_three_site_closure, def:mpdo_normalized_four_site_tail}
    Carried onto the tripartite site index, the three-site marginal of the
    normalized four-site chain is the three-site closure of $\mathcal K$
    against the normalized fourth-site tail $R_4$.
\end{theorem}

\begin{proof}
    \leanok
    \uses{thm:mpdo_reduced_four_three_tail}
    The entry formula (\ref{eq:rfp_three_site_marginal}) for
    $\sigma_3^{(4)}(\mathcal K)$, carried onto the tripartite site index
    $(i_1,i_2,i_3)$, gives
    \begin{align}
        \rho_{i_1i_2i_3,j_1j_2j_3}
        &=\tr\!\left(\mathcal K^{i_1j_1}\mathcal K^{i_2j_2}
          \mathcal K^{i_3j_3}R_4\right),
        \notag
    \end{align}
    which is the three-site closure condition
    (\ref{eq:rfp_three_site_closure}) against $R_4$.
\end{proof}

\begin{theorem}[Vanishing of the off-diagonal Hayashi sectors]
    \label{thm:mpdo_inverse_map_hayashi_sector_offdiagonal}
    \lean{MPOTensor.inverseMap_hayashi_sector_offdiagonal}
    \leanok
    \uses{def:mpdo_three_site_closure, def:mpdo_hayashi_inverse_factors}
    In the setting of
    Theorem~\ref{thm:mpdo_inverse_map_hayashi_sector_comparison}, for all
    $k\ne k'$, every off-diagonal sector satisfies
    \begin{align}
        R_{\beta_3,\alpha_1}
        \left[U_B\kappa_{\beta_1,\alpha_3}U_B^\dagger\right]_{(k,l,r),(k',l',r')}
        &=0.
        \label{eq:rfp_offdiagonal_vanishing}
    \end{align}
    Together with the diagonal comparison this gives the direct-sum structure
    of the factorized tensor: the direct sum over sectors is inherited from
    the splitting of the middle site.
\end{theorem}

\begin{proof}
    \leanok
    \uses{thm:mpdo_closure_outer_contraction}
    For $b$ in sector $k$ and $b'$ in sector $k'$ with $k\ne k'$, the Hayashi
    block equality at off-diagonal indices gives
    \begin{align}
        \sum_{i_2,j_2}(U_B)_{b,i_2}
        \rho_{i_1i_2i_3,j_1j_2j_3}
        \overline{(U_B)_{b',j_2}}
        &=0.
        \notag
    \end{align}
    Substituting this vanishing into the expansion of
    Theorem~\ref{thm:mpdo_closure_outer_contraction}
    gives (\ref{eq:rfp_offdiagonal_vanishing}).
\end{proof}

\begin{definition}[Sector tensors]
    \label{def:mpdo_sector_tensors}
    \lean{MPOTensor.sectorTensorL}
    \lean{MPOTensor.sectorTensorR}
    \leanok
    \uses{def:mpdo_hayashi_inverse_factors}
    Fix outer virtual indices $\alpha_1,\beta_3$ and a Hayashi sector $k$.
    The sector tensors are
    \begin{align}
        (l_k)_{\beta_1}
        &:=R_{\beta_3,\alpha_1}^{-1}\,p_k
          A^{(k)}_{\alpha_1,\beta_1},
        \notag\\
        (r_k)_{\alpha_3}
        &:=B^{(k)}_{\alpha_3,\beta_3}.
        \notag
    \end{align}
    % The source (arXiv:1606.00608, Appendix C.2, lines 1435--1448) leaves the
    % distribution of the scalar R^{-1} p_k between l_k and r_k unspecified;
    % here it is attached to l_k.
    These are the tensors of the factorization
    displayed in \cite[Appendix~C.2, lines~1435--1448]{Cirac2016MPDO_arXiv}.
\end{definition}

\begin{lemma}[Zero-weight left sector tensor]
    \label{lem:mpdo_sector_tensor_left_zero_of_weight_zero}
    \lean{MPOTensor.sectorTensorL_eq_zero_of_weight_eq_zero}
    \lean{MPOTensor.sectorEta_eq_zero_of_target_weight_eq_zero}
    \leanok
    \uses{def:mpdo_sector_tensors}
    If $p_k=0$, then $(l_k)_\beta=0$ for every virtual index $\beta$.
    Consequently, $\eta_{h,k}=0$ for every sector $h$.
\end{lemma}

\begin{proof}
    \leanok
    The scalar weight $p_k$ is a factor of every matrix entry of
    $(l_k)_\beta$. Each neighboring operator entering sector $k$ contains
    $(l_k)_\beta$ as its target factor.
\end{proof}

\begin{theorem}[Sector factorization of an injective simple tensor]
    \label{thm:mpdo_sector_factorization}
    \lean{MPOTensor.physicalSlice_sector_factorization}
    \leanok
    \uses{def:mpdo_sector_tensors}
    Suppose that $\mathcal K$ is injective, $\rho$ is its three-site closure
    against $R$, a Hayashi decomposition of $\rho$ has been chosen, and the
    outer indices satisfy $R_{\beta_3,\alpha_1}\ne0$. Then for all virtual
    indices $\beta_1,\alpha_3$,
    \begin{align}
        U_B\,\kappa_{\beta_1,\alpha_3}\,U_B^\dagger
        &=\bigoplus_k
          (l_k)_{\beta_1}\otimes(r_k)_{\alpha_3}.
        \label{eq:rfp_sector_factorization}
    \end{align}
    The direct sum refers to the physical indices and is inherited from the
    splitting of the middle site. This is the factorization displayed at
    \cite[Appendix~C.2, lines~1435--1448]{Cirac2016MPDO_arXiv}.
\end{theorem}

In the middle-site basis selected by $U_B$, the same factorization
(\ref{eq:rfp_sector_factorization}) is displayed as
\begin{center}
    % U_B kappa U_B^dagger = bigoplus_k l_k (x) r_k: U_B splits the
    % physical slice kappa_{beta_1,alpha_3}'s
    % physical space into the left sector factor l_k and right sector
    % factor r_k.  Each factor is an operator with one ket and one bra
    % leg.  Both sides have boundary signature
    % (virtual-west=1 | virtual-east=1 | physical-up=2 |
    % physical-down=2).
    \tenkzkernel
    \begin{tenkz}[rows={op,op,op}, cols=2]
        \tn[skin=pill, wide=2, ports={90@1:physical, 90@2:physical}]{U_B} \\
        \tn[skin=box, wide=2, ports={180:virtual, 0:virtual}]{\kappa_{\beta_1,\alpha_3}} \\
        \tn[skin=pill, wide=2, ports={270@1:physical, 270@2:physical}]{U_B^\dagger}
    \end{tenkz}
    $ = \bigoplus_k$
    \begin{tenkz}[cols=1, west=open]
        \tn[skin=box, ports={90:physical, 270:physical}]{l_k}
    \end{tenkz}
    $ \otimes $
    \begin{tenkz}[cols=1, east=open]
        \tn[skin=box, ports={90:physical, 270:physical}]{r_k}
    \end{tenkz}
\end{center}

\begin{proof}
    \leanok
    \uses{thm:mpdo_inverse_map_hayashi_sector_comparison,
        thm:mpdo_inverse_map_hayashi_sector_offdiagonal}
    On a diagonal sector, the sectorwise comparison gives
    $R_{\beta_3,\alpha_1}\widetilde\kappa^{(k)}
    =p_kA^{(k)}B^{(k)}$, and dividing by the nonzero entry
    $R_{\beta_3,\alpha_1}$ produces
    $(l_k)_{\beta_1}\otimes(r_k)_{\alpha_3}$.
    On an off-diagonal sector, the product with $R_{\beta_3,\alpha_1}$
    vanishes, and since $R_{\beta_3,\alpha_1}\ne0$ the sector itself
    vanishes.
\end{proof}

\begin{theorem}[Existence of the sector tensors]
    \label{thm:mpdo_exists_sector_factorization}
    \lean{MPOTensor.exists_physicalSlice_sector_factorization}
    \leanok
    \uses{def:mpdo_sector_tensors, def:mpdo_normalized_four_site_tail}
    Let $\mathcal K$ be injective with nonzero four-site trace, let $\rho$ be a
    three-site closure of $\mathcal K$ against the normalized fourth-site tail
    $R_4$, and choose a Hayashi decomposition of $\rho$. Then there exist
    sector tensors $l_k$ and $r_k$ such that, for all virtual indices
    $\beta_1,\alpha_3$,
    \begin{align}
        U_B\,\kappa_{\beta_1,\alpha_3}\,U_B^\dagger
        &=\bigoplus_k
          (l_k)_{\beta_1}\otimes(r_k)_{\alpha_3}.
        \notag
    \end{align}
    The neighboring operators $\eta_{k,h}$ are assembled from these tensors
    in Definition~\ref{def:mpdo_sector_eta}; the remaining steps of
    \cite[Appendix~C.2, Lemma~C.4]{Cirac2016MPDO_arXiv} prove the
    all-length identity and primitivity.
\end{theorem}

\begin{proof}
    \leanok
    \uses{thm:mpdo_sector_factorization,
        thm:mpdo_normalized_four_site_tail_entry_ne_zero}
    A nonzero four-site trace supplies outer indices with
    $(R_4)_{\beta_3,\alpha_1}\ne0$, and the sector factorization applies.
\end{proof}
